% Einheit 2 von 3 – Achsensymmetrie und Spiegeln (bank/symmetrie-abbildungen/e2.jsonl, zusammenbau v0.1)
%% TODO Zweigzeile Teil 1: was hier gelernt wird (Satz in Schülersprache aus dem Einheitstitel) – Einheit 2
\einheitenkopf[e2]{Einheit 2 von 3 · Achsensymmetrie und Spiegeln}
\zweigzeile{ab Kl. 5, je nach Buch bis Kl. 8 · P10}
\setcounter{aufgabe}{15}

%% TODO Titel: Ich-kann-Satz fehlt in der Bank (Platzhalter: Kettenname) – Nr. 16
%% TODO Anweisung: ein Satz über der ersten Teilaufgabe fehlt in der Bank – Nr. 16
\begin{aufgabe}{„Wie viele Achsen?“}
\begin{teile}
\teil Zeichne alle Symmetrieachsen der Figur mit dem Geodreieck ein. Zähle erst danach: Wie viele sind es?

\rechteck{5}{2}
\leerfeld Achsen
\end{teile}
\end{aufgabe}

%% TODO \verfahren{…}: Verfahrensüberschrift in Sachsprache fehlt in der Bank (2.3 b)
%% TODO Titel: Ich-kann-Satz fehlt in der Bank (Platzhalter: Kettenname) – Nr. 17
%% TODO Anweisung: ein Satz über der ersten Teilaufgabe fehlt in der Bank – Nr. 17
\begin{aufgabe}{Symmetrieachsen bestimmen}
%% TODO \rechenplatz{n}: Schrittzahl des Grundfalls fehlt in der Bank (nur bei fester Schreibform, 2.3 b)
\begin{teile}
\teil Passen beim Falten an der eingezeichneten Geraden beide Hälften genau aufeinander? \\ \kreuz{ja} \\ \kreuz{nein}

\drachen[achsen=senkrecht]{4}{3}{0.3}
\teil Zeichne die Symmetrieachse der Figur ein.

\drachen{4}{3}{0.3}
\teil Wie viele Symmetrieachsen hat die Figur?

\drachen{3.5}{2.5}{0.3}
\leerfeld Achse(n)
\teil Wie viele Symmetrieachsen hat die Raute?

\raute{5}{3}
\leerfeld Achsen
\teil Welche Figur hat keine Symmetrieachse? \\ \kreuz{Rechteck} \\ \kreuz{Parallelogramm} \\ \kreuz{Raute}
\teil Trage die Zahl der Symmetrieachsen ein. \\ \sachtabelle{lcccc}{Viereck & Quadrat & Rechteck & Raute & Drachenviereck}{Achsen & \leerzelle & \leerzelle & \leerzelle & \leerzelle}
\teil Wie viele Symmetrieachsen hat ein gleichseitiges Dreieck? \leerfeld Achsen
\teil Wie viele Symmetrieachsen hat der Großbuchstabe H? \leerfeld Achsen
\teil Im Drachenviereck ABCD ist die Diagonale BD die Symmetrieachse. Die Diagonale AC ist 6\,cm lang, F ist der Schnittpunkt der Diagonalen. Wie lang ist AF? AF = \leerfeld[cm]
\teil Die Platte eines Gruppentischs hat die Form eines gleichschenkligen Trapezes. Wie viele Symmetrieachsen hat sie? Kreuze an. \hfill (P10 2022 OS) \\ \kreuz{0} \\ \kreuz{1} \\ \kreuz{2} \\ \kreuz{3} \\ \kreuz{4} \\ \kreuz{5}

\trapez{6}{3}{2}
\teil Ein Kinderdrache hat die Form des Drachenvierecks ABCD. Seine Leisten sind die Diagonalen: AC ist 40\,cm, BD ist 70\,cm lang, sie kreuzen sich rechtwinklig. Wie viele Symmetrieachsen hat das Drachenviereck? Kreuze an. \hfill (P10 2025 OS) \\ \kreuz{0} \\ \kreuz{1} \\ \kreuz{2} \\ \kreuz{3} \\ \kreuz{4}

\drachen[diagonalen]{4}{2.3}{0.3}
\end{teile}
\end{aufgabe}

%% TODO \verfahren{…}: Verfahrensüberschrift in Sachsprache fehlt in der Bank (2.3 b)
%% TODO Titel: Ich-kann-Satz fehlt in der Bank (Platzhalter: Kettenname) – Nr. 18
%% TODO Anweisung: ein Satz über der ersten Teilaufgabe fehlt in der Bank – Nr. 18
\begin{aufgabe}{Spiegeln und Spiegelbild benennen}
%% TODO \rechenplatz{n}: Schrittzahl des Grundfalls fehlt in der Bank (nur bei fester Schreibform, 2.3 b)
\begin{teile}
\teil Punkt P liegt 3 Kästchen links von der Achse, sein Bildpunkt 3 Kästchen rechts, beide auf gleicher Höhe. Ist richtig gespiegelt? \\ \kreuz{ja} \\ \kreuz{nein}
\teil Spiegle den Punkt P an der Geraden a. Zähle die Kästchen bis zur Achse.

\begin{ksys}[xmin=0,xmax=9,ymin=0,ymax=8] \asymptote{x=4}{a} \punkt{2}{3}{P} \end{ksys}
\teil Spiegle das Dreieck ABC an der Geraden a.

\begin{ksys}[xmin=0,xmax=10,ymin=0,ymax=7] \asymptote{x=4}{a} \punkt{1}{1}{A} \punkt{3}{1}{B} \punkt{2}{4}{C} \end{ksys}
\teil Spiegle das Dreieck ABC an der Geraden a.

\begin{ksys}[xmin=0,xmax=8,ymin=-2,ymax=9] \gerade{0}{4}{a} \punkt{1}{1}{A} \punkt{4}{1}{B} \punkt{4}{3}{C} \end{ksys}
\teil Spiegle das Dreieck ABC an der Hochachse und gib die Koordinaten der Bildpunkte an.

\begin{ksys}[xmin=-6,xmax=6,ymin=-6,ymax=6] \punkt{2}{1}{A} \punkt{5}{1}{B} \punkt{4}{4}{C} \end{ksys}
A'( \leerfeld | \leerfeld ), B'( \leerfeld | \leerfeld ), C'( \leerfeld | \leerfeld )
\teil Die Punkte A, B, C, D sind die halbe Figur, a ist die Symmetrieachse. Ergänze die Figur zur achsensymmetrischen.

\begin{ksys}[xmin=0,xmax=8,ymin=0,ymax=7] \asymptote{x=4}{a} \punkt{4}{1}{A} \punkt{2}{1}{B} \punkt{2}{4}{C} \punkt{4}{6}{D} \end{ksys}
\teil Spiegle das Dreieck ABC an der schrägen Geraden a.

\begin{ksys}[xmin=0,xmax=8,ymin=0,ymax=8] \gerade{1}{0}{a} \punkt{1}{3}{A} \punkt{2}{5}{B} \punkt{1}{5}{C} \end{ksys}
\teil Ein rechtwinkliges Dreieck wird an einer seiner kurzen Seiten (Katheten) gespiegelt. Welche Figur bilden Dreieck und Spiegelbild zusammen?
\teil Nora schneidet ein Dreieck mit drei verschieden langen Seiten aus Papier aus. Sie legt es mit einer Seite an die Gerade a und spiegelt es an a. Welches Viereck bilden Dreieck und Spiegelbild zusammen? \hfill (P10 2018 OS)
\end{teile}
\end{aufgabe}

%% TODO Titel: Ich-kann-Satz fehlt in der Bank (Platzhalter: Kettenname) – Nr. 19
%% TODO Anweisung: ein Satz über der ersten Teilaufgabe fehlt in der Bank – Nr. 19
\begin{aufgabe}{prüfen, ob eine eingezeichnete Gerade Symmetrieachse ist (falten oder messen)}
\begin{teile}
\teil Ist eine der beiden eingezeichneten Diagonalen des Parallelogramms eine Symmetrieachse? Prüfe durch Messen. \\ \kreuz{ja} \\ \kreuz{nein}

\viereck[diagonalen]{(0,0)}{(4,0)}{(5.5,2.5)}{(1.5,2.5)}
\end{teile}
\end{aufgabe}

%% TODO Titel: Ich-kann-Satz fehlt in der Bank (Platzhalter: Kettenname) – Nr. 20
%% TODO Anweisung: ein Satz über der ersten Teilaufgabe fehlt in der Bank – Nr. 20
\begin{aufgabe}{Symmetrie der Vierecksarten im Haus der Vierecke zuordnen}
\begin{teile}
\teil Welche Vierecke im Haus der Vierecke haben mindestens zwei Symmetrieachsen?
\end{teile}
\end{aufgabe}

%% TODO Titel: Ich-kann-Satz fehlt in der Bank (Platzhalter: Kettenname) – Nr. 21
%% TODO Anweisung: ein Satz über der ersten Teilaufgabe fehlt in der Bank – Nr. 21
\begin{aufgabe}{Symmetrieachsen bestimmen – Fehler finden, Begründen, Darstellungswechsel, Anwendung}
\begin{teile}
\teil Finde den Fehler. Sven zeichnet in ein gleichschenkliges Trapez zwei Symmetrieachsen ein: eine senkrechte und eine waagerechte durch die Mitte, wie beim Rechteck.
\teil Begründe, warum ein Parallelogramm keine Symmetrieachse hat, obwohl es gleichmäßig aussieht.
\teil Zeichne ein Viereck mit genau einer Symmetrieachse und zeichne die Achse ein.

\begin{ksys}[xmin=0,xmax=8,ymin=0,ymax=6] \end{ksys}
\teil Ein Logo für die Schülerzeitung soll genau zwei Symmetrieachsen haben. Welche Grundform passt? \\ \kreuz{Quadrat} \\ \kreuz{Rechteck} \\ \kreuz{gleichseitiges Dreieck}
\end{teile}
\end{aufgabe}
